\section{A chamber with no marked graph filling}
\label{ch:section}

The obstruction will first be placed in a compact chamber.  Its boundary
marking is essential: the obstruction does not merely assert that some
disk problem has no solution.

\begin{proposition}[The marked chamber]
\label{ch:chamber}
There are a compact connected oriented smooth four-manifold $D$, its
connected boundary $S$, a finitely generated free group $H$, and a map
$c:D\to BH$ with the following properties.
\begin{enumerate}
\item The map $c$ induces an isomorphism on fundamental groups, and $D$
has the homotopy type of a finite graph with $2k$ two-spheres attached
at its basepoint.  Its module $\pi_2D$ has a specified hyperbolic basis
over $\Z H$, represented by framed immersed spheres with vanishing
quadratic self-intersection obstructions.
\item Attach a three-cell to each basis sphere, and denote the resulting
space by $P$.  Then $(P,S)$ is an oriented Poincar\'e pair of formal
dimension four, and $P\simeq BH$.  The boundary marking is
$c|_S:S\to BH$.
\item There is no compact oriented topological four-manifold $D'$ with
boundary identified with $S$ and a homotopy equivalence $c':D'\to BH$
whose restriction to $S$ is homotopic to $c|_S$.
\end{enumerate}
\end{proposition}

Here and below, sphere classes and intersections use fixed whiskers and
deck translations.  A hyperbolic pair has zero squares, mutual product
$1$, and zero framing and Wall self-intersection obstructions.  We prove
the proposition without using an embedded surgery on $D$.

\subsection{The geometric and tensor inputs}
\label{ch:inputs}

We use the marked geometric construction of \cite[Section~2]{TO}, rather
than only its main disk-nonembedding theorem. The marked bodies and
exterior duals are supplied by its Lemmas~2.3 and~2.4. Its
Proposition~2.2 and Lemmas~2.5--2.6 give the cut consequences under
hypothetical disk replacements; those conditional conclusions are
distinct from the marked geometric construction itself.
We recall the precise data needed here.
Fix $m\geq5$ and $p\geq9$.  Start with a four-dimensional one-handlebody
on $2m+p$ handles, attach cancelling two-handles to the first $2m$
handles with product framings, and plumb these two-handles in $m$ pairs,
interchanging the two coordinate directions at every plumbing.
The resulting smooth oriented manifold $X$ has graph homotopy type.
Its cover associated to the plumbing quotient $\pi_1X\to F_m$ is a
tree of blocks, each a one-handlebody on $p$ extra generators; neighboring
blocks meet in a product four-ball.  The exposed boundary of a block is
\[
 K_v=\#^p(S^1\times S^2)\setminus
       \operatorname{int}\nu(L_{2m}),
\]
where $L_{2m}$ is an unlink in a ball.  Its port meridians
$x_1,\ldots,x_{2m}$ and extra loops $l_1,\ldots,l_p$ form a free basis.
At each join the meridian on one side is the longitude on the other.

Write $\mathcal C=\{1,l_1,\ldots,l_p\}$.  The marked-body construction
installs a disjoint two-stage grope body for every ordered triple of
distinct ports and every independent choice
\[
 X_a=c_a x_{r_a}^{\delta_a}c_a^{-1},\qquad
 c_a\in\mathcal C,\quad \delta_a\in\{1,-1\},\quad a=1,2,3.
\]
Its lower surface is a punctured torus: one branch has a single cap,
and the other carries the upper punctured torus with its two caps.
All tip tracks, conjugating paths, and product framings are specified.
The initial boundary word $R$ satisfies
\begin{equation}
\label{ch:laws}
 R|_{X_a=1}=1\quad(a=1,2,3),\qquad
 R\equiv[X_1,[X_2,X_3]]\pmod{\Gamma_4F(X_1,X_2,X_3)}.
\end{equation}
Internal changes of the marked paths are allowed by these identities;
one must retain the entire finite family, not replace it by one
unconjugated commutator.

Let $F_1,\ldots,F_k$ be the lower punctured tori, with symplectic bases
$(a_j,b_j)$, where $a_j$ is the single-cap branch.  Delete their open
product tubes, including the proper edge collars, to obtain $E$.
Lemma~2.4 (Algebraic duals in the exterior) of \cite{TO} supplies
the truncated single caps $f_j$ and framed spheres $g_j$ in $E$ such
that
\begin{equation}
\label{ch:exterior-duals}
 \lambda(f_i,g_j)=\delta_{ij},\qquad
 \lambda(g_i,g_j)=0,\qquad \widetilde\mu(g_j)=0
 \quad\text{over }\Z[\pi_1E].
\end{equation}
The framing on $f_j$ is the surface-compression framing.  The $g_j$
are compressed normal-circle tori over the $b_j$ curves; their
individual representatives have trivial normal bundle.  Equation
\eqref{ch:exterior-duals} is a group-ring statement before any
meridians are killed.

We also recall exactly how the cut obstruction is used.  Suppose that a
compact smooth oriented four-manifold with the marked graph type of
the parity cover of $X$ has two connected vertex pieces $A,B$, with
outside faces $K_A,K_B$, and $2m$ connected three-dimensional cuts
$Y_i$, each with its prescribed coordinate torus as its entire boundary.
Orient $Y_i$ as a face of $A$, so its orientation as a face of $B$
is opposite.
Suppose the differences of restrictions give an isomorphism in degree
two and, in degree one, an isomorphism after the extra-loop evaluations
are set to zero.  Suppose further that precisely $m$ joins are active:
at an active join the surviving coordinate meridian is at $A$, whereas
the other coordinate is the meridian at $B$.  Require every marked law
\eqref{ch:laws} in its respective vertex group, and the free ambient
background in which the port meridians are trivial and the named extra
letters can be assigned independently.

These data cannot exist.  Indeed, Proposition~6.1 (Vertex data) of
\cite{TO} gives a characteristic-zero field $K$ (one may take
$\mathbb C((\sigma))$), separate coordinate blocks $h_i$ of total
dimension $N$, and algebraic power-series germs $F_i(h_i)$ and
$b_i(h_i)$, with $b_i$ indexed by the active set $I$, $|I|=m$.
The germs are normalized by $F_i(0)=b_i(0)=0$.  Set
\begin{equation}
\label{ch:tensor-ring}
 T=K[t_i:i\in I]/(t_i^2:i\in I),\qquad
 S_0=\sum_iF_i,\qquad S_t=S_0+\sum_{i\in I}t_i b_i.
\end{equation}
There are formal parametrizations
\[
 H_L(x,t)\in(x,t)T[[x]]^N,\qquad
 H_R(y)\in(y)K[[y]]^N,
\]
where $x$ has $d$ coordinates and $y$ has $N-d$ coordinates.  Their
central tangent maps $\partial_xH_L(0,0)$ and $\partial_yH_R(0)$
are injective with complementary images in $K^N$; nilpotent constant
displacements of $H_L$ are allowed.  They satisfy
\begin{equation}
\label{ch:tensor-vanishing}
 S_t(H_L)=0,\qquad S_0(H_R)=0.
\end{equation}
For every three distinct active indices $i,j,k$, they satisfy
\begin{align}
 t_it_jt_k&\in
 \bigl((\partial_{h_\nu}S_t)(H_L):\nu=1,\ldots,N\bigr)
                  \subset T[[x]],\label{ch:tensor-left}\\
 (b_ib_jb_k)(H_R)&\in
 \bigl((\partial_{h_\nu}S_0)(H_R):\nu=1,\ldots,N\bigr)
                  \subset K[[y]].\label{ch:tensor-right}
\end{align}
Theorem~3.1 (Finite tensor obstruction) of \cite{TO} rules out
\eqref{ch:tensor-ring}--\eqref{ch:tensor-right} for $m\geq5$.
In particular, the memberships are in the actual gradient ideals,
not their radicals, and $T$ retains all mixed square-free monomials.
The independent conjugators in the full law list are what supply all
these triple identities. These two source results, together with the
marked edge potentials, algebraicity and formal comparison of
\cite[Proposition~4.5, Theorem~5.1 and Proposition~5.2]{TO}, are the
tensor obstruction used here.

\subsection{The exterior with its side marking}
\label{ch:exterior}

The imported geometry supplies the surfaces and their algebraic duals.
We now use those data to construct $D$, keeping track of the actual
tube-side maps. This marking will identify the boundary of a hypothetical
filling with the loop data needed for the tensor contradiction.

We make one placement refinement of the marked-body construction.
Place its longitudinal tip annuli on disjoint parallels outside
sufficiently thin attaching solid tori in the initial one-handlebody
boundary, and put the finitely many connectors in their complement.
This is possible by the marked handlebody model: the bodies are boundary
joins of the separated annular collars.  The complement of the attaching
solid tori survives in the boundary after the two-handles are attached;
the plumbing ports are on the new belt-handle portions.  Thus the
unpushed lower surfaces lie in the final exposed boundary of $X$.
The independent ribbon twists matching the cap framings are made in
the separated parallel tubes.  They do not pass a body across a
plumbing join.  The returning tip collars alone are lengthened back
to product parallels on the attaching regions and then into the cores.
The construction and intersection calculation of
\eqref{ch:exterior-duals} are unchanged.

Consequently the $F_j$ are proper inward pushes of disjoint boundary
surfaces, with returning collars at their initial boundary curves.
Moreover $F_j\to X$ is null-homotopic.  The three tip loops bound the
corresponding core caps in $X$; hence the bottom $a_j$ and $b_j$ loops
die in $\pi_1X$.  Since $X$ has graph homotopy type and $F_j$ retracts
to a graph, this proves the assertion about the whole surface map.

\begin{lemma}[The marked exterior model]
\label{ch:exterior-model}
With the surface-product trivializations, the exterior and its tube-side
maps have the homotopy-pushout model
\[
 E\simeq X\cup_{\coprod_jF_j}\coprod_j(F_j\times S^1),
\]
where each gluing map to $F_j\times S^1$ is the inward section.
The displayed $F_j\times S^1$ maps represent the actual tube sides.
\end{lemma}

\begin{proof}
It is enough to describe one surface; the constructions have disjoint
collars.  Write $B=\partial X$ and $K=\partial F$.  Before rounding
corners, put the pushed surface at a horizontal collar level, with
vertical walls along $K$ returning to the boundary.  The inner part of
the exterior has homotopy type $X$, the outer part has homotopy type
$B\setminus K$, and their overlap has homotopy type $B\setminus F$.
Product slices, thickened by a variable positive width off the deleted
surface, give this homotopy-pushout decomposition.

In $B\setminus K$, take a two-sided neighborhood of
$F^\circ=F\setminus K$ whose width decreases to zero near $K$.
Relative to $B\setminus F$, it is a bridge over $F^\circ$, joining
the two sides of the cut.  Replace $F^\circ$ in this diagram by a
slightly shrunken copy $F_0$; this is an objectwise homotopy equivalence.
The actual deleted surface and $K$ are unchanged.  In the rectangular
normal coordinates to its horizontal copy, split the tube circle over
$F_0$ into its outer and inward semicircles.  The outer semicircle is
the bridge.  The inward semicircle lies in the inner $X$ piece and
gives a specified homotopy between the bridge's endpoint maps.

Using that homotopy to identify the endpoints replaces the bridge by
$F_0\times S^1$, attached along its inward section.  Its circle is
precisely the outer bridge followed by the inward return, hence the
actual normal circuit.  Finally $F_0\to F$ and
$F_0\times S^1\to F\times S^1$ are homotopy equivalences.  This
also accounts for the proper edge collars and proves the assertion
about the side maps, not just the absolute homotopy type of $E$.
\end{proof}

For each $j$, regard a solid torus $V_j$ as a compression body from
the punctured torus $F_j$ to a disk, compressing $a_j$.  Define
\[
 D=E\cup_{\coprod_j(F_j\times S^1)}
                  \coprod_j(V_j\times S^1).
\]
All gluings use the just specified product framings; round corners.
The remaining disk face of $\partial V_j$ contributes a solid torus
to the new boundary.  Hence $S=\partial D$ is the surface-framed
surgery of $\partial X$ on the initial grope-boundary link.  In
particular it is connected.

By Lemma~\ref{ch:exterior-model},
\[
 D\simeq X\cup_{\coprod_jF_j}\coprod_j(V_j\times S^1).
\]
Each $F_j\to X$ is null-homotopic.  The extra summand is therefore
the homotopy cofiber of $F_j\to V_j\times S^1$.  Retract the latter
space to the torus on its longitude $b_j$ and normal meridian $t_j$.
The cofiber adds one two-cell killing $b_j$ and one on the null loop
$a_j$.  Cancel the $b_j$ one/two-cell pair.  The torus two-cell then
has null attaching map, so the result is
\[
 S^1_{t_j}\vee S^2\vee S^2.
\]
Thus $H=\pi_1D$ is the original free graph group with one new free
meridian $t_j$ per surface, and $D$ has the graph-and-spheres homotopy
type in Proposition~\ref{ch:chamber}.

\subsection{The middle form and the Poincar\'e pair}
\label{ch:duality}

Close $f_j$ by a meridian disk of $V_j$ at its fixed circle coordinate,
obtaining a sphere $u_j$.  The two compression framings agree, so its
second Stiefel--Whitney evaluation is zero.  The $g_j$ remain in $E$;
therefore \eqref{ch:exterior-duals} gives
\[
 \lambda(u_i,g_j)=\delta_{ij},\qquad \lambda(g_i,g_j)=0
 \quad\text{over }\Z H.
\]
The hermitian matrix of the $u_i$ has even diagonal.  As $H$ has no
element of order two, every such diagonal is $z+\bar z$: choose one
representative from each nonidentity inverse pair and halve the even
identity coefficient.  Add suitable combinations of the $g_j$ to the
$u_i$, taking a triangular half off the diagonal and this half on it.
The resulting spheres, denoted $a_i$, together with $b_i=g_i$, have
hyperbolic intersection matrix.  The usual framed representative
adjustments give zero Wall self-intersections: choose normal Euler
number zero using the evenness, and then use
$\lambda(a_i,a_i)=\mu(a_i)+\overline{\mu(a_i)}$ and the absence of
order-two elements.  This is the ordinary immersed-sphere calculus
of \cite{FQ}; no disk-embedding assertion enters it.

These $2k$ classes form a basis.  Indeed, $\pi_2D$ is already known
to be free of rank $2k$.  Its square matrix of the displayed classes
has a left inverse obtained from their nonsingular pairing.  To see
directly that it also has a right inverse, pass to each finite quotient
of $H$.  The resulting matrices act on finite-dimensional complex
vector spaces, where a left inverse is a right inverse.  Residual
finiteness of the free group separates the finite support of any
nonzero group-ring element, so the right-inverse identity holds over
$\Z H$ itself.  We henceforth use this specified basis.

Attach a three-cell on every $a_i$ and $b_i$ to form $P$.  A map from
the graph with $2k$ two-spheres to $D$ realizing this basis is a
homotopy equivalence: it induces an isomorphism on $\pi_1$ and on the
homology of simply connected covers.  Thus $P\simeq BH$.

We check Poincar\'e duality for the pair, since this step must not be
confused with embedded surgery.  Set $\Lambda=\Z H$.  In the duality
of $(D,S)$, the middle sphere module maps isomorphically onto a split
summand of $H_2(D,S;\Lambda)$, by its nonsingular form.  Attaching the
three-cells quotients both absolute and relative homology by these
middle summands; the attaching map is injective on them, so no new
third homology appears.  On cohomology the pullbacks identify the new
groups with the annihilators of the sphere summands.  If
$f:(D,S)\to(P,S)$ is the natural cell-attachment inclusion and
$[P,S]=f_*[D,S]$, cap-product naturality reads
\[
 z\frown[P,S]=f_*\bigl(f^*z\frown[D,S]\bigr).
\]
Hence the old duality isomorphisms induce exactly the new ones on
these annihilators and quotients.  The spaces have finite free chain
complexes over $\Lambda$, so the regular-coefficient cap-product
criterion gives Poincar\'e duality with arbitrary local coefficients.
This proves parts (1) and (2) of Proposition~\ref{ch:chamber}.

\begin{corollary}
\label{ch:boundary-homology}
The map $\pi_1S\to H$ is onto.  If $r=\operatorname{rank}H$, then
$H_1(S;\Z)\to H_1(P;\Z)$ is an isomorphism and
$H_2(S;\Z)\cong\Z^r$.
\end{corollary}

\begin{proof}
Duality and the graph type of $P$ give
$H_1(P,S;\Z H)=H_0(P,S;\Z H)=0$.  Therefore
$H_0(S;\Z H)\to H_0(P;\Z H)$ is an isomorphism.
The former is free abelian on the cosets of the image of $\pi_1S$,
and the latter is $\Z$; there can be only one coset.
With ordinary coefficients the pair sequence and duality give
$H_2(P,S)=H_1(P,S)=0$ and
$H_3(P,S)\cong H^1(P)\cong\Z^r$.  The assertions follow.
\end{proof}

\subsection{Topological fillings and the cut obstruction}
\label{ch:forbidden-filling}

The construction and duality calculation establish the first two parts
of Proposition~\ref{ch:chamber}. For the third part we must analyze an
arbitrary topological filling with the specified boundary map. The
following cut construction supplies the hypotheses recalled in
Section~\ref{ch:inputs}; it does not presume a smooth structure on that
filling.

Suppose a marked filling $D'$ as in part (3) existed.  Attach the
dual two-handles that undo the surgery on the initial link.  Denote
the result by $W$, so $\partial W=\partial X$.  Each attaching circle
represents its prescribed new free generator $t_j$ under the marking
$D'\simeq BH$.  The corresponding one/two-cell pair cancels in the
homotopy type.  Thus $W$ has the graph type of $X$, with the original
boundary-to-graph map: equality of the marked fundamental-group maps
gives equality up to homotopy because the graph is aspherical.
The two-handle cocores are disjoint proper locally flat disks bounded
by the original initial link.  In particular all the original marked
words are now represented by disk boundaries.

Here is the cut construction with the hypotheses on $W$ made explicit.
Use the marked quotient $\pi_1W\to F_m$ and its associated tree cover.
Prescribe the equivariant map to be constant at the appropriate vertex
on each outside face $K_v$ and each lifted disk neighborhood, and
to cross the appropriate edge on each boundary torus collar.  These
prescriptions agree where they meet.  A section of the associated
bundle with contractible tree fiber extends them to the rest of $W$.
The free boundary extra loops, with whiskers, generate the fundamental
group of this tree cover, just as for the marked graph model of $X$.

We must allow $W$ to be topological.  By Quinn's punctured smoothing
theorem \cite[Corollary~2.2.3]{QuinnEnds}, delete an interior point
away from the disks and extend the given boundary smoothing to the
complement.  Prescribe the map to be constant near this point before
extending the section.  Approximate it smoothly over the open middle
intervals of the finitely many graph edges, away from the disks and
puncture, and choose regular values.  The cuts are compact smooth
three-manifolds, disjoint from the disks.  Their filled complementary
pieces are topological four-manifolds, smooth away from the finite
set of lifted punctures in each finite quotient.

For completeness, connectedness is not automatic from a regular-level
construction.  Form the component graph of the cuts in the tree cover.
Its fundamental group is a quotient of that of the cover, by joining
successive crossings through connected complementary pieces.  Every
extra boundary generator lies in one piece, so maps trivially to this
graph.  It follows that the component graph is a tree.  The outside
faces and torus collars give an embedded principal copy of the original
plumbing tree: different vertex labels lie in different components of
the complement of the target midpoints.  Any remaining branch has a
unique nearest principal vertex.  Its stabilizer is trivial, since it
fixes that vertex.  The quotient has finitely many components; hence
these branches have finite valence and uniformly bounded depth, and
each is finite.  Ignore their closed cuts and merge them into their
principal vertices.  This leaves connected vertex pieces and one
connected cut with each prescribed torus as its entire boundary.
The disk prescriptions have not changed.  This is precisely the
argument of Lemma~2.5 (Connected tree cuts) of \cite{TO}, now
applied to the marked $W$ rather than requiring a homeomorphism $W=X$.

Pass to the parity quotient, sending every plumbing generator to the
nonidentity element of $\Z/2$.  Its filled ambient space $W_2$ has
the homotopy type of the graph with two vertices and $2m$ edges,
and $p$ extra circles at each vertex.  Write $A,B$ for the filled
vertex pieces and $Y_1,\ldots,Y_{2m}$ for their cuts.  Over any
characteristic-zero field $K$, Mayer--Vietoris gives
\begin{equation}
\label{ch:restriction-two}
 H^2(A;K)\oplus H^2(B;K)
       \xrightarrow{\ \cong\ }\bigoplus_iH^2(Y_i;K),
\end{equation}
because $H^2(W_2;K)=H^3(W_2;K)=0$.  In degree one the difference
of restrictions is onto.  Its kernel is the image of $H^1(W_2;K)$
modulo the $2m-1$ graph classes, and evaluation on the $p$ named
extra loops at each vertex identifies that kernel with $K^{2p}$.
Thus the restriction map becomes an isomorphism on the subspace
where these extra evaluations vanish.

On a cut $Y_i$, half-lives/half-dies says that the image of
$H_1(\partial Y_i;K)$ in $H_1(Y_i;K)$ is a line.  If both original
coordinate loops had nonzero image, their images would be nonzero
multiples.  A class from either vertex annihilates the longitude
bounding in its outside face, and hence would annihilate both
coordinates.  The degree-one surjectivity just proved would then
force every class on $Y_i$ to do so, a contradiction.  Hence one
original coordinate $\beta_i$ dies and the other $\alpha_i$ survives.
The plus endpoint is the side at which $\beta_i$ is a longitude.
The parity involution pairs the two lifts of each join and exchanges
their plus endpoints, giving exactly $m$ active edges at $A$.

The prescribed disks lie in their own vertex pieces.  They therefore
give every law \eqref{ch:laws} at those vertices, with all the
independently chosen conjugators and signs.  The marked ambient free
group still has trivial port meridians and freely assignable named
extra letters.  In particular it supplies the common flat background
used in the Vertex data Proposition: assign eight extra letters the
eight nonidentity matrices among
$\exp(a\sigma e)\exp(b\sigma f)$, $a,b\in\{0,1,2\}$, for the
standard nilpotent generators $e,f$ of $\mathfrak{sl}_2$.
Together with the identity, their adjoint operators span all
endomorphisms over $\mathbb C((\sigma))$, as in
\cite[Proposition~6.1]{TO}.
Thus the geometric and marked hypotheses of the tensor interface
in Section~\ref{ch:inputs} have all been established.

It remains to justify that its vertex calculation works for the
punctured smooth structures just used.  We give this step rather
than assuming that the filled vertices admit smooth structures.

\begin{lemma}[Forms on a punctured vertex]
\label{ch:puncture-forms}
Let $V$ be one of the filled compact vertex pieces and let $Z$ be
its finite set of punctures.  In degree zero take smooth functions
on $V\setminus Z$ separately constant near each puncture, and in
positive degrees take smooth forms vanishing near the punctures.
This differential graded algebra computes $H^*(V;\mathbb C)$,
naturally on the unpunctured faces.  It supports the vertex
calculation in the Vertex data Proposition of \cite{TO}, with the
filled cohomology groups and with the same Stokes identities.
\end{lemma}

\begin{proof}
Regard these forms as a sheaf on $V$.  Its stalk at a puncture is
$\mathbb C$ in degree zero and zero in higher degrees; elsewhere
it is the ordinary de Rham resolution.  Partitions of unity can
be chosen constant on smaller neighborhoods of the finitely many
punctures.  The complex is therefore a fine resolution of the
constant sheaf, so computes the filled cohomology.  A based condition
in degree zero removes only $H^0$.  Linear splittings onto cohomology
and contracting homotopies may consequently be chosen inside this
complex, exactly as for the ordinary vertex de Rham complex.

Wedges, Lie brackets, the curvature recursions, and differentiation
in the formal parameters preserve this algebra.  Each fixed
positive-degree coefficient vanishes near $Z$, so integration and
Stokes take place on a compact smooth subdomain containing its
support.  The additional end boundaries contribute zero.  An entire
formal series need not have one common vanishing neighborhood:
every calculation at a fixed coefficient uses only finitely many
terms.  The source's formal lifts and divisions likewise use
finite linear combinations coefficientwise, and extensions from
face collars can be supported away from $Z$.

Removing interior points in dimension four does not change the
fundamental group.  At each finite nilpotent order the flat systems
in question vanish near the punctures in their local trivializations;
gauges between them can be separately constant at those ends.
Thus the based holonomy and exact-gauge arguments use the same
marked group as the filled vertex.  Allowing separate constants
at distinct ends is important here.  These are the uses of vertex
smoothness in the proof of Proposition~6.1 of \cite{TO}.
Its edge-algebraicity theorem, \cite[Theorem~5.1]{TO},
is still applied only to the unchanged compact smooth
three-manifolds $Y_i$.
\end{proof}

Apply Lemma~\ref{ch:puncture-forms} to the two vertex pieces.
Equation~\eqref{ch:restriction-two} and the degree-one calculations
refer to their filled cohomology, which the lemma computes.  The
Vertex data Proposition, \cite[Proposition~6.1]{TO}, now produces exactly
\eqref{ch:tensor-ring}--\eqref{ch:tensor-right}.  The Finite tensor
obstruction contradicts their existence.  Therefore $D'$ cannot
exist, completing the proof of Proposition~\ref{ch:chamber}.
